\section{Limitations}\label{sec:limitations}

\paragraph{One synthetic generator family.} All three tiers come from a single procedural vol-and-jump generator; no real price series enters the environment, and the tiers differ in that generator's parameters, not in kind. F4 quantifies the cost: the tape has fat tails on the stressed tiers but weak volatility clustering at episode length, and the finite-panel-calibrated three-check conjunction passes on only 1 of 24 seeded panels. The Fano ratio is conditional-IID-calibrated but exploratory and does not gate. The existing Calm calibration grid produces no preset that meets its diagnostic, confirmation and volatility rule. An agent that excels here has demonstrated skill against this generator's structure, not against a market, and the realism gate says so rather than hiding it. \Cref{sec:validation} argues that a gate a committed public generator can fail is evidence the gate discriminates, which is a claim about the instrument and not a defence of the tape: on the four-level empirical-validity ladder of \citet{fagiolo2019validation} the canonical generator sits at Level 0 on the gated axis, and every finding built on it is instrument calibration on uncertified tape. The paragraph below states which experiments run on other processes instead.

\paragraph{The interface guarantee is not an information guarantee.} The leak-freedom of \cref{sec:leakfree} is a statement about the interface's shape: no operation returns a post-cursor bar. It is not a claim that public deterministic scenarios are unpredictable. \Cref{sec:predictability} makes the distinction empirical. A causal AR(5) adversary stays near its unconditional baseline, but an exhaustive scan of a public $2^{16}$ seed band recovers all 16 tested seeds from one observed bar in about one second. Full-space brute force remains infeasible. A complete 64-bit SplitMix64 output inverts to one state, whereas its published 53-bit unit representation leaves $2^{11}$ candidates; tape-level inversion through the later price arithmetic remains open. The opt-in sealed derivation of \cref{sec:sealed-seeds} defeats this specific bounded-band scanner 0/16 when its salt is withheld, and a forward protocol can commit the salt hash before entry. The evidence here is a simulated workflow, not evidence that any live evaluation used it. Any public bounded-band evaluation should be treated as compromised by this scanner until it documents sealed seed custody, a pre-entry commitment, non-reuse and a reveal deadline. The derivation is PRF-style, not a certified MAC, so its security rests on salt entropy and secrecy rather than a cryptographic proof.

\paragraph{The baselines are trivial policies, not trained agents.} Every number in \cref{sec:validation,sec:findings} is produced by fixed reference policies (buy-and-hold analogs, one-step tilts, closed-form allocators) and two deliberately biased behavioral counterparties. They fix the zero point of each instrument and calibrate its resolution; they do not establish what a trained agent scores, and this paper's claims are scoped to environment and protocol for that reason. The constraint is the learner's, not the simulator's: \cref{sec:protocol-compute} measures the environment at 14{,}240 steps per second on the slowest of the three shipped surfaces, so a PPO run is bounded by the policy's own training cost rather than by environment time. In particular, the near-zero within-tier generalization gaps in F3 are the expected control for policies with nothing to overfit, not evidence that the environment resists overfitting by learners. Relatedly, no reference policy is rank-eligible under the kernel's full conjunction (\cref{sec:f1}). \Cref{sec:witness} establishes that the eligibility set is non-empty on every tier, but it does so with an out-of-band oracle that replays the seed one bar ahead, which is not an entrant and is not scored on any board; no policy that reads only the observation interface has yet been shown eligible.

\paragraph{Named future experiments.} Two follow-ups remain specific. (1) A learning baseline: PPO trained through the shipped Gymnasium adapter on the train band, evaluated on a named disjoint evaluation band and cross-regime, so F3 contains a policy that can actually overfit. (2) Tape-level analytic state recovery from quantized SplitMix64-derived prices. The exact-finalizer inverse and the resulting $2^{11}$ candidates per 53-bit unit are tested, but no observed price has been inverted; sealed seeds (\cref{sec:sealed-seeds}) defeat this specific bounded-band scanner but do not answer this one. The manipulation positive control in \cref{sec:f5-positive-control} finds a profitable sampled round trip at $\beta = 0.5$ and none among its sampled linear-impact cells, though the search did not cover all interactions or transient impact. The rank-eligibility existence question of \cref{sec:f1} is answered by the witness of \cref{sec:witness}, by an oracle rather than by an entrant.

\paragraph{The Python probe layer is outside the golden-hash guarantee.} The byte-identical claim covers the Rust core, the WebAssembly build and the generated scenario bytes returned through the Python binding: the native and Python test suites assert all seven committed golden fingerprints and the WebAssembly suite asserts the canonical and clustered ones, and direct native--WebAssembly parity tests enforce the second. The Python binding delegates environment dynamics to that same native Rust core and has deterministic replay tests, but the golden coverage on the Python side is the scenario serialization, not the package's other serialized artifacts. The Python reductions that produce the realism, manipulation, adverse-selection, failure, ecology and predictability evidence are deterministic in their seeds, but are off the golden-hash path; their floating-point reproducibility is the platform's, not the paper's. The concave-impact engine itself is also excluded from the canonical byte-identity claim because a general exponent uses \texttt{powf}; the linear default remains on the untouched golden path.

\paragraph{The realism diagnostic does not certify every experiment's process.} F4 applies to the canonical position-trading generator. The Avellaneda--Stoikov market-making process in F2 and the market-clearing models in F5 and F6 are separate simulators with their own structural and positive-control checks; they are not passed through F4. Their results are therefore calibrations of those specified models, not claims conditioned on the canonical tape's realism verdict. The distinction matters in both directions: F4's failure does not mechanically invalidate a different process, and the separate process is not market-valid merely because its internal controls pass.

\paragraph{The impact models are stylized.} The endogenous market composes Kyle and Almgren--Chriss terms; the adverse-selection module does not separate a meta-order's information from its impact; the follower population is a one-parameter caricature. The concave ablation is not a universal test: its dimensionless flow $|Q/V|$ sits two to three orders of magnitude below the unit crossover, so exponents below one amplify impact relative to the linear arm at the tested calibration, and the symmetric schedule does not search the asymmetric round trips nonlinear theory permits. \Cref{sec:f5-positive-control} does search them and finds cells that pay at $\beta = 0.5$ with no followers, so the symmetric null says this pump does not pay here, and the positive control says the instrument can see schedules that do; the extended sweep enlarges the sampled axes but does not cover interactions, other flow scales or transient impact.

\paragraph{The guard is a deny-list.} The structural guarantee covers the environment's own surface. \texttt{LookaheadGuard} extends it to harness tools by pattern, and a deny-list is exactly as good as its patterns; a novel leak path through a side channel the environment does not own (wall-clock time, filesystem state, a network call) is out of scope for both layers. The replay check verifies decisions against frozen data and therefore catches falsified results, not information an agent obtained elsewhere.

\paragraph{No live external entrants yet.} The contract, the conformance kit and the replay verifier exist, and the evidence run exercises them, but no third party has yet submitted an agent. A hosted intake, a live leaderboard and multi-tenant isolation of untrusted agent processes do not exist, and the governance document is a written promise backed by tests that has not been stressed by a hostile or careless external implementer. The significance claim of this paper is therefore the verifiability infrastructure that is checkable today; adoption is future work, not an assumption.

\paragraph{Broader impact.} Two dual-use surfaces deserve an explicit norm. First, leaderboard scores: a byte-verifiable score is more persuasive, not more externally valid, and a verified SharpeArena number is a claim about a synthetic environment only. SharpeArena scores are not performance marketing, and presenting them to retail audiences as evidence of real-market skill would be a misuse the project's documentation states as out of bounds. Second, the manipulation probe: it is a payoff-search harness over manipulation schedules, which is dual-use by construction. We judge publishing it net-positive because the linear-impact arm is a diagnostic against the precise no-manipulation assumptions stated in \cref{sec:f5}, not a search for deployable trades; all 45 sampled linear configurations lose money, but the paper makes no guarantee outside that finite grid. The module carries its simulator-diagnostic disclaimer in code, and the crude, legible schedules it implements are the ones surveillance already catches in real markets, so the probe teaches defenders more than attackers. The profitable concave cells use sliced accumulation and block liquidation, a schedule family motivated by the published manipulation literature rather than a claim about an executable market strategy.
